\uuid{1DwU}
\exo7id{2318}
\titre{exo7 2318}
\auteur{matexo1}
\organisation{exo7}
\datecreate{2002-02-01}
\isIndication{false}
\isCorrection{false}
\chapitre{Calcul d'intégrales}
\sousChapitre{Changement de variables}
\cpge{mpsi-mp2i.integration.c, pcsi.integration.b, ptsi.integration.b, mpsi-mp2i.calcul-algebrique-trigonometrie.d, pcsi.calcul-algebrique-trigonometrie.d, ptsi.calcul-algebrique-trigonometrie.d, bcpst1.derivees-primitives-integrales.d, bcpst1.trigonometrie}
\module{Analyse}
\niveau{L1}
\difficulte{}

\contenu{
\question{
Sans calculer les int\'egrales, montrer que
$$ \int_0^{\pi /2} \sin^n x dx = \int_0^{\pi /2} \cos^n x dx. $$
}
}
